La enseñanza de la trigonometría en educación media arrastra una tensión que no logra resolverse. El currículo pide que el estudiantado razone sobre variación, periodicidad y articulación entre registros. En el aula ocurre otra cosa: predominan las fórmulas, los procedimientos y los ejercicios cerrados \parencite{diaz-fernandezEnsenanzaTrigonometria4o2014,rosaEnsenanzaAprendizajeTrigonometria2023}.

La consecuencia es previsible, pero no por eso menos preocupante. Los estudiantes resuelven triángulos rectángulos sin mayor dificultad. Interpretar una gráfica trigonométrica ya es otro asunto. ¿Por qué se rompe ahí la comprensión? Trabajar razones trigonométricas no parece bastar para construir la comprensión funcional que la educación media espera.

La literatura ofrece una pista: estas dificultades se relacionan, en parte, con una enseñanza centrada en el cálculo mecánico y con poco espacio para argumentar \parencite{fiallolealAcercaInvestigacionEducacion2015,delgadoEstrategiaDidacticaPara2023}. Tratar seno, coseno y tangente como funciones reales —y no solo como razones que se calculan— exige tareas distintas. Exige también discusiones más profundas y criterios de evaluación pensados para la comprensión, no solo para el resultado \parencite{torres-corralesResignificacionRazonTrigonometrica2021,escalantegodoyUsoComprensivoRazones2018}.

Hay errores que se repiten con la insistencia de un patrón: confundir razón trigonométrica con función, entender mal la periodicidad, no lograr conectar representación gráfica, expresión algebraica y valor numérico \parencite{zubietaTipificacionErroresDificultades2018,ruedaAproximacionEnsenanzaRazones2012,iriarteDificultadesAprendizajeEcuaciones2025}. En Colombia los reportes de evaluación llevan años documentando estas mismas brechas \parencite{ministeriodeeducacionnacionalmenInformeNacionalSaber2019,icfesProgramaParaEvaluacion2024}.

¿El problema está solo en el estudiante? No parece razonable pensarlo así. El tipo de tarea que se propone también pesa. Cuando una actividad exige explorar, comparar y justificar, se abre espacio para comprender. Cuando se limita a repetir, el aprendizaje se reduce a ejecutar procedimientos \parencite{benitezAprendizajeTrigonometriaMediante2022,delgadoAulaInvertidaRendimiento2024,zambrano-choezEstrategiasParaEnsenanza2025}.

La tecnología, aquí, ofrece algo real. GeoGebra permite ver el cambio en tiempo real y conecta representaciones que de otro modo quedan aisladas \parencite{calderonTecnologiasDigitalesPara2024,agudelocarvajalFortalecimientoProcesoEnsenanza2020}. La IA podría ampliar ese potencial: retroalimentación oportuna, rutas de exploración distintas según quien aprende. Pero solo si la tarea conserva exigencia matemática. Reducida a seguir instrucciones, su aporte al pensamiento trigonométrico es mínimo \parencite{romeroalonsoRolInteligenciaArtificial2024,castillodelpezoUseArtificialIntelligence2024}.

No hay, todavía, evidencia propia de cómo se sostiene ese equilibrio —entre apoyo de la IA y exigencia matemática— en un aula real. Ese vacío motiva este trabajo.

De estas premisas nace la pregunta que define el problema: ¿de qué manera el diseño e implementación de tareas trigonométricas integradas con Inteligencia Artificial puede favorecer el paso desde el cálculo hacia el desarrollo del pensamiento trigonométrico en estudiantes de media?